\documentclass[11pt]{amsart}
\usepackage[utf8]{inputenc}
\usepackage{amsmath,amssymb,amsthm}
\usepackage[margin=1in]{geometry}
\newtheorem{proposition}{Proposition}
\DeclareMathOperator{\Oz}{Oz}
\DeclareMathOperator{\Aut}{Aut}
\DeclareMathOperator{\Ad}{Ad}
\newcommand{\kk}{\Bbbk}
\title{Exact ozone computations for three type (14641) algebras at $h=1$}
\author{James G\'omez \and Helbert Venegas}
\date{November 22, 2026}
\begin{document}
\maketitle

% Supplementary note — James Gomez and Helbert Venegas
% Independent R family presentation: Zhang--Zhang (2009), family R;
% checked against https://4d-as.ncag.info/families/ore-r/ on 2026-11-22.
% Q and V use the Gomez--Venegas draft, Table tab:Mt.
\section{Three ozone computations for the specializations $h=1$}
\noindent Throughout $\kk$ is algebraically closed of characteristic zero. We use
$\Oz(B)=\Aut_{Z(B)\text{-alg}}(B)$ and the right-normal convention
$w a=\sigma(a)w$, so $\sigma=\Ad_w$ in the quotient division ring.
The following statements are \emph{specializations}, not generic-twist theorems.

\begin{proposition}
For each of $B=\mathbb Q(1),\mathbb R(1),\mathbb V(1)$,
$\Oz(B)\cong C_4$.
\end{proposition}
\begin{proof}
For $\mathbb Q(1)$, use the separate proof in \texttt{QRV\_H1\_THEOREMS.tex}:
$b=x_1(y_1^2+y_2^2-2y_1y_2)+x_2(y_1^2-y_2^2)$ is central;
the full quadratic automorphism equations together with its fixedness leave
precisely the four powers of $\Ad_{x_1}$.

For $\mathbb V(1)$ let
\[
 w=x_1(y_1^2-y_2^2)-x_2(y_1^2+y_1y_2).
\]
Direct PBW reduction, from the four mixed relations of the normal-element
localization of the supplied manuscript, gives
\[
 w x_1=(-x_1+x_2)w,\quad w x_2=x_2w,\quad
 w y_1=-y_2w,\quad w y_2=y_1w.
\]
Hence $\sigma=\Ad_w$ restricts to
\[
 X=\begin{pmatrix}-1&1\\0&1\end{pmatrix},\qquad
 Y=\begin{pmatrix}0&-1\\1&0\end{pmatrix},
\]
and has order four. The existing upper-bound calculation is as follows:
any ozone automorphism preserves both degree-one planes; since $x_2$,
$v=x_1^2(x_1-x_2)^2$, and $s_y=y_1^2+y_2^2$ are central,
its $x$-matrix belongs to
$\{I,\left(\begin{smallmatrix}-1&1\\0&1\end{smallmatrix}\right)\}$
and its $y$-matrix is a signed permutation. The six defining quadratic
relations restrict this finite candidate set to four elements. Therefore
$|\Oz(\mathbb V(1))|\le4$, and the powers of $\Ad_w$ realize equality.

For $\mathbb R(1)$, in generators $x_1,x_2,y_1,y_2$,
\[
 x_2x_1=-x_1x_2,\quad y_2y_1=-y_1y_2,
\]
\[
\begin{array}{ll}
 y_1x_1=x_1y_1+x_2y_1+x_1y_2,& y_1x_2=x_1y_2,\\
 y_2x_1=x_2y_1,&y_2x_2=-x_2y_1-x_1y_2+x_2y_2.
\end{array}
\]
The PBW reducer certifies centrality of
\[
 s_x=x_1^2+x_2^2,\quad s_y=y_1^2+y_2^2,\quad
 z=x_1y_1+x_1y_2+x_2y_1-x_2y_2.
\]
Any ozone automorphism preserves each degree-one plane, and the
$(-1)$-quantum-plane relations imply that its two matrices are monomial.
Fixing $s_x$ and $s_y$ restricts all their nonzero entries to $\pm1$.
Of the 64 pairs, only eight preserve the four mixed relations; fixing $z$
leaves four. Define
\[
 u=x_2y_1+x_1y_2.
\]
A second exact PBW reduction gives
\[
 u x_1=x_2u,\quad u x_2=-x_1u,\quad
 u y_1=-y_2u,\quad u y_2=y_1u.
\]
As $u\ne0$ is normal, $\Ad_u\in\Oz(\mathbb R(1))$. It has order four;
its four distinct powers exhaust the above candidate set.
\end{proof}

\paragraph{Scope.} The $\mathbb R$ presentation was transcribed from
\texttt{https://4d-as.ncag.info/families/ore-r/} and should be
cross-checked against the original Zhang--Zhang matrix $\Sigma$ before
submission. Exact source verification is distinct from exact symbolic
computation. None of these arguments establishes $\Oz(\mathbb Q(h))$,
$\Oz(\mathbb R(h))$ or $\Oz(\mathbb V(h))$ for arbitrary PI twists.

\end{document}
